\chapter{Conclusioni}
\label{chap:conclusioni}

\section{Il lavoro svolto}
\label{sec:riepilogo}

L'obiettivo del lavoro era riuscire a dotare Alchemist della capacità
di esporre ai nodi di una simulazione grandezze geospaziali misurate,
in funzione della posizione in cui essi si trovano e dell'istante di
simulazione in cui le richiedono, ottenendole dai data store Copernicus.
Il risultato ottenuto è un nuovo modulo del simulatore, organizzato
nei quattro sottomoduli descritti nel \Cref{chap:progettazione} e
dichiarabile interamente nella configurazione della simulazione.

Il sotto-obiettivo di rappresentazione è stato soddisfatto da \texttt{RasterGrid}
e \texttt{GridSnapshots}, che descrivono rispettivamente una griglia di campioni
e una successione di griglie omogenee, indipendenti dal formato e dalla fonte
da cui tali dati provengono. Queste astrazioni portano con sé l'estensione
spaziale e temporale nelle quali la grandezza di riferimento è nota, e trattano
in modo esplicito le celle prive di misura (\req{F7}); delle griglie sono fornite
una realizzazione densa e una sparsa, scelte dinamicamente durante la costruzione
della serie in funzione della densità delle singole griglie.

Il sotto-obiettivo di risoluzione è stato soddisfatto dal layer e dalle strategie
a cui esso delega. Il layer stabilisce alla costruzione la corrispondenza fra
il tempo di simulazione e gli istanti del mondo reale, dichiarata da chi
configura la simulazione anziché fissata dal modulo (\req{F3}),
trova i campioni pertinenti all'istante in cui viene interrogato e incarica
la produzione del valore a una strategia di interpolazione e a un convertitore,
entrambi intercambiabili e dichiarabili dalla configurazione (\req{F5}, \req{F6}, \req{F11}).

Il sotto-obiettivo di acquisizione è stato soddisfatto dal sottosistema
costituito da \texttt{CacheKey}, \texttt{ExternalDataProvider} e \texttt{CacheManager},
e dalle loro realizzazioni che dialogano con i tre data store Copernicus.
L'identità stessa di una richiesta ne decide deterministicamente la voce
di cache (\req{NF4}), cosa che rende possibile riusare dati già acquisiti
senza alcuna richiesta di rete (\req{F10}); tutta l'acquisizione avviene
al caricamento della simulazione, in modo tale che ogni condizione che
impedisca di disporre dei dati si manifesti prima che l'esecuzione cominci (\req{NF6}).
Un risultato aggiuntivo riguarda il protocollo dei data store Copernicus:
il loro comportamento aderisce solo in parte a \ac{OGC} API
Processes, e gli scostamenti riscontrati, discussi nel \Cref{chap:implementazione},
non sono documentati dal servizio: sono stati verificati empiricamente, a partire
dalle risposte reali raccolte dallo sviluppo, e la loro descrizione
resta valida anche per altri scenari d'utilizzo.

Il lavoro descritto ha potuto realizzarsi come aggiunta
al simulatore, senza modifiche alle
astrazioni del suo modello (\req{NF1}).

La validazione copre tutte e tre le parti: una suite di test
che non dipende dai data store né dal possesso di credenziali (\req{NF3}),
la verifica del caricamento da configurazione, l'uso del client sui data
store reali, una simulazione dimostrativa per confermare il corretto comportamento
del sistema di cache e la corretta lettura ed esposizione dei valori.

\section{I limiti della soluzione}
\label{sec:limiti}

Il modulo opera entro alcune ipotesi e alcune scelte prese.

\begin{itemize}
  \item \textbf{La forma della griglia}: l'ipotesi stabilita nella
  \Cref{subsec:griglie-supportate} esclude dataset campionati su griglia
  nativa in proiezione cartografica. Tali dati restano utilizzabili solo se
  ricondotti esternamente a una griglia di assi indipendenti e caricati poi
  tramite il secondo costruttore del layer.

  \item \textbf{Le dimensioni del dato}: la grandezza esposta deve
  dipendere dalle sole tre dimensioni tempo, latitudine e longitudine. I dataset
  che sono definiti anche su altre dimensioni, come quota o livello di pressione,
  non possono individuare un valore unico per una posizione e un istante, e restano
  quindi fuori dal campo di applicabilità.

  \item \textbf{L'occupazione di memoria}: la realizzazione fornita della serie
  di griglie costruisce l'intera successione in memoria al caricamento,
  scelta che risponde a \req{NF7} ma il cui costo cresce con l'estensione temporale
  e geografica della richiesta.

  \item \textbf{La sola lettura}: il valore esposto cambia soltanto perché il
  tempo di simulazione avanza lungo una serie di griglie già disponibile. Il dato
  sottostante non può quindi essere aggiornato durante l'esecuzione.
  Resta quindi fuori portata ogni
  scenario in cui l'ambiente conserva l'effetto di ciò che i nodi hanno fatto,
  ad esempio una simulazione in cui dei depuratori riducono la concentrazione
  degli inquinanti nella zona in cui operano.

  \item \textbf{La conoscenza preventiva del nome della variabile}: quando i file
  contengono più di una grandezza, la configurazione deve indicare quale far esporre
  al layer, e il nome da indicare è quello interno al file, che non coincide
  con quello utilizzato per fare la richiesta verso il data store.
  Nella modalità che acquisisce i dati da un data store tale nome è noto prima
  che l'acquisizione sia avvenuta: chi conduce la simulazione deve eseguirla una
  prima volta, o ispezionare i file ottenuti, per apprenderlo. Il caso è più marcato
  con i file GRIB, nei quali la grandezza è identificata da codici numerici e il
  nome esposto dipende dalle tabelle che la libreria di lettura possiede
  (\Cref{sec:formati-raster}).
\end{itemize}

\section{Sviluppi futuri}
\label{sec:sviluppi-futuri}

La struttura del modulo rende possibile alcune direzioni di sviluppo future senza
modificarne le parti già esistenti.

La prima è il superamento della restrizione sulla forma della griglia.
Rappresentare dataset in proiezione cartografica richiede di invertire la proiezione
fra il sistema di riferimento del dataset e quello geografico, in modo da poter
individuare la cella contenente una determinata posizione; ciò si traduce
in una nuova realizzazione di \texttt{RasterGrid}, che il resto del modulo
può usare tramite la stessa astrazione.

La seconda riguarda le fonti di dati. Il sottosistema di acquisizione
conosce i data store Copernicus solo nelle realizzazioni che ne descrivono
una richiesta e li interrogano: utilizzare una fonte diversa richiede unicamente
di implementare la coppia \texttt{CacheKey} ed \texttt{ExternalDataProvider}
corrispondente. Il caso più immediato è il recupero autonomo da parte del simulatore
degli estratti cartografici utili ad alimentare un \texttt{OSMEnvironment}, che oggi
devono essere predisposti manualmente.

Una terza direzione è una realizzazione della serie di griglie che carichi
i dati su richiesta invece che per intero all'avvio della simulazione,
con un compromesso tra occupazione di memoria e costo di lettura.
Il layer non richiederebbe alcuna modifica poiché basterebbe dare una nuova
implementazione di \texttt{GridSnapshots}.

Le ultime due dipendono invece da interventi sul simulatore.
Una revisione del motore di simulazione renderebbe possibile l'impiego
di layer aggiornabili durante l'esecuzione. Una migrazione del tipo di tempo
di Alchemist (della quale la \Cref{subsec:corrispondenza-temporale} riporta lo studio di fattibilità),
o un sistema di conversione generica dal tempo di simulazione
a una rappresentazione arbitraria del tempo, semplificherebbe la corrispondenza
fra i due domini temporali incontrati nel lavoro svolto; poiché tale 
corrispondenza è stabilita unicamente nel layer, l'impatto di eventuali
modifiche sul dominio temporale resterebbero circoscritte a esso.
